% =============================================================
%  CHAPTER: Methodology
%  Thesis: Hallucination Detection in Large Language Models
%          via Modern Hopfield Retrieval over SwiGLU Memory Banks
%  Status: Revised. Primary bank is the gate matrix A^{(l)} = W_1^{(l)T}
%          (consistent with fundamentals.tex, Eq. memory_bank_definition);
%          retrieval scores are a fixed per-row rescaling of the gate
%          pre-activations (derived view, s != g).
% =============================================================

\label{chap:methods}

This chapter operationalises the retrieval framework of
Chapter~\ref{chap:theory}: it specifies how the energy, entropy, and
divergence quantities defined there are extracted from a frozen model and
assembled into a hallucination score.

The method is developed within four bounding conditions: factoid question
answering with short gold answers; decoder-only models whose feed-forward
sublayers use the SwiGLU parameterisation (Section~\ref{subsec:ffn}); English
datasets; and hallucinations that appear in the generated answer rather than
in upstream chain-of-thought tokens. These conditions define the regime in
which the proposed signal is claimed to apply. Generalisation to other tasks
is out of scope.

Six further constraints govern the design itself. First, the weights and
activations of the model under study are fully accessible. Second, no
fine-tuning of the model is permitted. Third, the Hopfield features are
computed from a single prefill pass followed by a single greedy decoding
pass, although sampling-based baselines may use multiple forward passes.
Fourth, the detection mechanism is a read-only observer: the decoding
distribution must be identical with and without detection, so no
intervention on generation is allowed. Fifth, each feature must admit a
per-layer interpretation traceable to a specific weight matrix or hidden
state, in keeping with the mechanistic interpretability objective. Sixth,
the pipeline must be deterministic for fixed random seeds, model
checkpoints, and chat-template configurations, within a fixed precision and
software stack on fixed hardware.

% =============================================================
\section{Problem Description}
\label{sec:problem_description}
% =============================================================

Let $f_\theta$ be a frozen, instruction-tuned, decoder-only language model
with $L$ layers and SwiGLU feed-forward sublayers
(Section~\ref{subsec:ffn}). Given a question $q$ and a set of gold
short-form answers $\mathcal{A}_q = \{a_1, \ldots, a_m\}$ from a factoid
question-answering dataset, the model produces an answer
$\hat{y} = f_\theta(q)$ by greedy autoregressive decoding, selecting the
highest-probability token at each step. The question is wrapped in a fixed
chat template, a structured prompt format that prepends role markers and a
system turn so that $q$ reaches $f_\theta$ in the same conversational form
used during instruction tuning.

A \emph{hallucination} is a binary indicator $z_q \in \{0, 1\}$, with
$z_q = 1$ when $\hat{y}$ is not factually supported by $\mathcal{A}_q$ and
$z_q = 0$ otherwise. These labels are needed to evaluate the scoring
function defined below, and obtaining them at scale by hand is the main
bottleneck for any large evaluation corpus. The label is therefore assigned
by a separate language model acting as an external judge, queried through a
locked prompt template that asks whether $\hat{y}$ is supported by
$\mathcal{A}_q$. The model under study is never asked to evaluate its own output.

Because the judge is queried through an external service whose outputs are
not deterministic, the labelling step is reproducible only given the
archived label set rather than bit-for-bit on re-execution. Label quality
is guarded by validating judge agreement against a human-labelled subset,
reported as Cohen's $\kappa$ in Chapter~\ref{chap:results} together
with the subset size.

The detection task is to construct a scoring function
\begin{equation}
    s : (q,\,\hat{y},\,\mathcal{I}(f_\theta, q, \hat{y}))
    \,\longmapsto\, \mathbb{R},
    \label{eq:detection_function}
\end{equation}
where $\mathcal{I}(\cdot)$ is a fixed family of internal observables
collected from the forward passes of $f_\theta$ on $q$ and during the
generation of $\hat{y}$. Higher values of $s$ indicate a higher predicted
probability that $z_q = 1$, and the quality of $s$ is measured by the area
under the receiver operating characteristic curve (AUROC) against the
labels $z_q$.

% =============================================================
% =============================================================
% =============================================================
% =============================================================
\section{Problem Analysis}
\label{sec:problem_analysis}
% =============================================================

The detection function of Equation~\eqref{eq:detection_function} is
underdetermined as stated: it does not specify which internal observables
enter $\mathcal{I}$, how they are compared across the two phases of
inference, or how the comparison is reduced to a scalar score. The task
decomposes into five sub-problems, and their dependencies run in sequence.
The choice of probe positions determines what is extracted; the memory-bank
definition determines what those extractions are compared against; the
retrieval scoring converts both into energies and probability distributions;
the divergence quantification compares those distributions across the
prefill and generation phases; and the aggregation and probing step
assembles the per-layer quantities into the scoring function $s$. This
section analyses each sub-problem in turn, identifying what makes it
nontrivial and what requirements any resolution must satisfy. The decisions
that satisfy these requirements are stated in
Section~\ref{sec:solution_design}, and their computational realisation is
deferred to Section~\ref{sec:implementation}. Following the convention of
Section~\ref{sec:hopfield}, single hidden states are treated as column
vectors in this chapter, so retrieval scores appear as
$\mathbf{M}^{(l)}\mathbf{r}_t^{(l)}$; entrywise these are the same inner
products as the row-form expression
$\mathbf{r}_t^{(l)}{\mathbf{M}^{(l)}}^{\!\top}$ of
Chapter~\ref{chap:theory}.

% ----------------------------------------------------------------
\subsection{Choice of Probe Positions}
\label{subsec:pa_probe_positions}
% ----------------------------------------------------------------

The retrieval framework of Chapter~\ref{chap:theory} assigns a query to
every layer and every token position, but the detection hypothesis is
comparative: it concerns the relationship between the retrieval structure
established while the model processes the grounded question and the
structure induced while it produces the answer. Autoregressive inference
supplies this contrast naturally. During prefill, the forward pass over the
prompt, the input tokens are externally given and grounded by construction;
during generation, each forward pass conditions on tokens the model itself
has emitted, so the input becomes progressively self-referential. Any
operationalisation of the hypothesis must therefore designate, per layer, a
reference object drawn from the prefill and a trajectory of objects drawn
from the generation steps.

The generation side admits little freedom: the quantities of interest exist
once per newly produced token, and discarding positions would discard
signal. The prefill side is genuinely underdetermined. The prompt contains
many token positions, each inducing its own retrieval distribution, while
the comparison requires a single per-layer reference. A pooled reference
over the question span would average distributions induced by tokens of
heterogeneous grammatical and informational roles, producing a mixture that
no single forward-pass state ever realises, and the comparison against
per-token generation states would then be one between objects of different
kinds. A single-position reference avoids this mismatch, and among the
prompt positions only the final one conditions on the complete question; it
is also the state from which the first answer token is decoded, so it
represents the retrieval configuration at the moment the model commits to
its answer. The requirement is thus a single-token, last-position prefill
reference, with the further condition that prefill and generation states be
extracted under an identical prompt format, since any discrepancy in the
conditioning context would contaminate the phase comparison with a prompt
artefact.

% ----------------------------------------------------------------
\subsection{Identifying the Memory Bank}
\label{subsec:pa_memory_bank}
% ----------------------------------------------------------------

The second sub-problem is inherited directly from
Section~\ref{subsec:gated_ffn_memory}. In the ReLU architecture analysed by
\citet{geva2021transformer}, the first weight matrix is the unique
structural candidate for the stored keys, because each intermediate
coefficient is a function of a single inner product with one of its rows. In
the gated architecture of Equation~\eqref{eq:gated_ffn_generic}, each
coefficient depends on two projections of the input simultaneously, so no
weight matrix occupies that position by necessity and the bank must be fixed
by argument. The analysis of Chapter~\ref{chap:theory} establishes the
grounds: the gate branch alone passes through the nonlinearity and performs
the selective activation that defines a key, while the value branch
modulates magnitude and the down projection writes content to the residual
stream. Whatever bank is adopted must respect this division of roles, must
be a fixed, non-trainable function of the frozen weights so that the
read-only constraint of this chapter is preserved, and must yield one stored
pattern per intermediate unit so that the retrieval distribution has a
per-unit interpretation.

This sub-problem also fixes the scope of what the detector can claim to
observe. A bank built from the recognition step measures which stored
patterns a token engages, not what the layer subsequently writes back to the
residual stream; the latter is a separate quantity, and any conclusion drawn
from the method must be phrased at the level of pattern engagement. The
hypothesis under test is correspondingly delimited: grounded and
hallucinated generation are conjectured to differ in this pattern of
engagement, measurable as a shift in the retrieval distribution between the
prefill reference and the generated tokens.

% ----------------------------------------------------------------
\subsection{Scale Confounds in Retrieval Scoring}
\label{subsec:pa_scale_confounds}
% ----------------------------------------------------------------

Converting queries and bank into energies and distributions exposes the
method to two scale confounds, and each imposes a requirement on the scoring
construction.

The first concerns the query. The Modern Hopfield energy of
Equation~\eqref{eq:modern_hopfield_energy} contains the quadratic term
$\tfrac{1}{2}\|\mathbf{q}\|_2^2$, so any systematic difference in
hidden-state norm between the phases would shift the energy even if the
alignment structure were unchanged. A detector built on raw energies could
then succeed merely as a norm-change detector, a degenerate explanation that
the analysis must be able to exclude. Two requirements follow. The query
must be taken at a point where its scale is controlled; the probe point
$\mathbf{r}_t^{(l)}$ of Equation~\eqref{eq:probe_point} satisfies this,
since RMSNorm fixes the root-mean-square scale up to the learned vector
$\boldsymbol{\gamma}$ and makes queries comparable across tokens and phases.
And the energy must remain attributable to its components, so the scoring
construction has to keep the log-sum-exp term and the norm term separable
rather than reporting only their sum. At the same time, the query norm is
not pure nuisance: the energy depends on it explicitly, and a change in
representation magnitude between phases is itself a candidate signal.
Whether the query should be normalised is therefore a genuine design
decision, trading the removal of a confound against the loss of an
informative degree of freedom, and it must be resolved explicitly rather
than inherited from convention.

The second concerns the stored patterns. Bank rows are learned weight
vectors with heterogeneous norms, so a raw inner product confounds the
direction of a pattern with its magnitude, and large-norm rows dominate the
softmax of Equation~\eqref{eq:retrieval_distribution} regardless of
directional match. Comparability across memory slots requires placing the
rows on a common scale. Doing so has a consequence that must be stated
wherever faithfulness to the model is claimed: once the rows are rescaled,
the resulting scores are no longer identical to the pre-activations the
model itself computes, but a per-slot rescaling of them. The construction
remains a deterministic function of the model's native quantities, yet it is
a derived view rather than a raw activation, and the formal statements of
the method must keep this distinction explicit.

A related decision point is where along the gate branch the score should be
read. The nonlinearity transforms the inner products before they enter the
remainder of the layer, but the Hopfield energy and update of
Section~\ref{sec:hopfield} are defined on the inner products themselves;
reading the score after the nonlinearity would substitute a transformed
quantity into a formalism built for the linear one and would sever the
correspondence on which the energy interpretation rests. The requirement is
that the scoring preserve the dot-product form assumed by the framework.

% ----------------------------------------------------------------
\subsection{Comparing Distributions Across Phases}
\label{subsec:pa_distribution_comparison}
% ----------------------------------------------------------------

The fourth sub-problem is the choice of comparison quantities, and the
constraints of the setting narrow it considerably. The comparison should not
depend on which phase is designated as reference, which rules out asymmetric
measures as primary features; the Kullback--Leibler divergence of
Section~\ref{subsec:kl_divergence} is accordingly suited as a building block
rather than a feature. Features feeding a downstream classifier should be
bounded, so that a handful of extreme values cannot dominate the fit, and
they should remain numerically stable where probabilities are small, since
softmax outputs over banks with $K_l$ in the tens of thousands place most
mass on few entries. No single divergence captures every aspect of
distributional change, and a probe trained on one measure inherits its blind
spots, so the feature set should combine measures with complementary
sensitivity; the properties established in Section~\ref{sec:divergences}
identify the Jensen--Shannon divergence and the squared Hellinger distance
as a pair satisfying these conditions. Finally, single-distribution
summaries face a comparability problem of their own: the bank size $K_l$
varies across layers and architectures, so raw entropies live on different
scales, and only the normalised entropy of
Equation~\eqref{eq:normalised_entropy} supports per-layer and cross-model
comparison.

% ----------------------------------------------------------------
\subsection{From Layerwise Quantities to a Score}
\label{subsec:pa_aggregation}
% ----------------------------------------------------------------

The final sub-problem is the reduction of per-token, per-layer quantities to
the scalar score $s$, and its central analytical point is that the direction
of the effect cannot be assumed. It is tempting to posit that hallucination
corresponds to degraded retrieval, that is, to higher energy and more
diffuse distributions during generation. The fixed-point structure of
Section~\ref{subsec:classical_to_modern_hopfield} does not support this as
an assumption: a hallucinated answer may equally correspond to a low-energy,
sharply concentrated retrieval in which the query has settled into the basin
of a well-formed but factually wrong association, so that the failure is one
of content rather than of retrieval stability. Furthermore, there is no a
priori reason for the sign or strength of the effect to be uniform across
layers, given the depth-dependent specialisation of feed-forward keys
reported by \citet{geva2021transformer}. The mapping from features to score
must therefore be learned from labelled data, free to assign either
direction to each layer, while remaining simple enough that any
discrimination it achieves is attributable to the features themselves rather
than to representational capacity in the classifier; this preserves the
per-layer interpretability constraint stated at the opening of this chapter.
Aggregation over the generated tokens must additionally not conflate answer
length with signal strength, which requires a length-invariant reduction of
the per-token quantities, and features of heterogeneous scales must be
placed on a common scale before fitting, since the residual query-norm
dependence identified in Section~\ref{subsec:pa_scale_confounds} leaves
energies on layer-dependent scales.

Each requirement assembled in this section has a one-to-one resolution in
the design that follows: a last-token prefill reference compared against
per-token generation states under an identical prompt format; a fixed,
recognition-side memory bank derived from the frozen weights; a scoring
construction that preserves the dot-product form, controls pattern scale,
makes the derived-view caveat explicit, and keeps the energy decomposition
separable; bounded, symmetric divergence pairs with normalised entropy; and
a learned, direction-agnostic linear probe on standardised, length-invariant
features. Section~\ref{sec:solution_design} states these decisions formally.
% =============================================================
% =============================================================
\section{Solution Design}
\label{sec:solution_design}
% =============================================================

The five sub-problems of Section~\ref{sec:problem_analysis} now combine into
a concrete detector. This section defines the retrieval scoring
construction, the per-layer feature set, the calibration of the inverse
temperature, the detection probe with its evaluation protocol, and the full
algorithmic pipeline.
% ----------------------------------------------------------------
\subsection{Retrieval Scoring Construction}
\label{subsec:retrieval_scoring}
% ----------------------------------------------------------------

The scoring construction resolves the requirements identified in
Section~\ref{subsec:pa_scale_confounds}: it must preserve the dot-product
form assumed by the Hopfield framework, control the scale of the stored
patterns, resolve the query-normalisation decision explicitly, and make
explicit that the resulting scores are a derived view of the model's native
activations rather than those activations themselves.

The model computes, at every layer $l$ and every token position $t$, the
raw gate pre-activations
\begin{equation}
    \mathbf{g}_t^{(l)}
    \;=\;
    {\mathbf{W}_1^{(l)}}^{\!\top}\mathbf{r}_t^{(l)}
    \;=\;
    \mathbf{A}^{(l)}\mathbf{r}_t^{(l)}
    \;\in\; \mathbb{R}^{K_l},
    \label{eq:gate_preactivations}
\end{equation}
where $\mathbf{r}_t^{(l)}$ is the probe point of
Equation~\eqref{eq:probe_point} and $K_l$ is the intermediate dimension of
layer $l$. These are the inner products the forward pass passes to the Swish
nonlinearity; reading them before that transformation preserves the
dot-product form on which the Hopfield energy and update rule are defined.

The retrieval scores against the memory bank $\mathbf{M}^{(l)}$ of
Equation~\eqref{eq:memory_bank_definition} are obtained by rescaling each entry of
$\mathbf{g}_t^{(l)}$ by the fixed L2 norm of the corresponding bank row,
\begin{equation}
    s_{t,i}^{(l)}
    \;=\;
    \frac{g_{t,i}^{(l)}}{\|\mathbf{m}_i^{(l)}\|_2}
    \;=\;
    {\hat{\mathbf{m}}_i^{(l)}}{}^{\!\top}\mathbf{r}_t^{(l)},
    \qquad
    \hat{\mathbf{m}}_i^{(l)} \coloneqq
    \frac{\mathbf{m}_i^{(l)}}{\|\mathbf{m}_i^{(l)}\|_2},
    \qquad i = 1, \ldots, K_l,
    \label{eq:retrieval_scores}
\end{equation}
so that $s_{t,i}^{(l)} = \|\mathbf{r}_t^{(l)}\|_2\cos\theta_i$, where
$\theta_i$ is the angle between $\mathbf{r}_t^{(l)}$ and the $i$-th gate
row. The vector $\mathbf{s}_t^{(l)} = (s_{t,1}^{(l)}, \ldots,
s_{t,K_l}^{(l)})^\top \in \mathbb{R}^{K_l}$ collects all scores at
position $(l,t)$. The rescaling is query-independent and fixed at bank
construction; the forward pass is never modified. The query itself is not
normalised: the factor $\|\mathbf{r}_t^{(l)}\|_2$ is retained as an
informative component of the score, resolving the design decision posed in
Section~\ref{subsec:pa_scale_confounds}, while the separable energy
decomposition keeps its contribution attributable.

Two consequences of this construction must be stated explicitly. First, the
scores $\mathbf{s}_t^{(l)}$ are not identical to the gate pre-activations
$\mathbf{g}_t^{(l)}$ the model computes: they are a per-slot rescaling of
them by norms that are fixed but not generally equal to one. The
construction is a deterministic function of the model's native quantities
and preserves the read-only constraint, yet it is a derived view rather than
a raw activation. Any claim of faithfulness to the model must be qualified
accordingly. Second, the Hopfield energy and retrieval distribution of
Equations~\eqref{eq:modern_hopfield_energy}--\eqref{eq:retrieval_distribution}
are evaluated by the detector at the pair
$(\mathbf{M}^{(l)},\,\mathbf{r}_t^{(l)})$ using $\mathbf{s}_t^{(l)}$ as
the score vector; no update iteration is performed, and no claim is made
that the model itself executes Hopfield retrieval dynamics.
% ----------------------------------------------------------------
% ----------------------------------------------------------------
\subsection{Per-Layer Feature Design}
\label{subsec:feature_design}
% ----------------------------------------------------------------

For each layer $l \in \{0, \ldots, L-1\}$ and each sample $q$, let
$\mathbf{r}^{(l,\mathrm{ref})}$ be the prefill last-token hidden state and
let $\{\mathbf{r}_t^{(l)}\}_{t=1}^{T_g}$ be the hidden states over the
$T_g$ generated tokens, with $T_g \le T_{\max}$ a fixed cap on the number
of tokens the model may produce.  Let $E^{(l,t)}$, $H^{(l,t)}$, and
$\bar{H}^{(l,t)}$ be the Hopfield energy
(Equation~\eqref{eq:modern_hopfield_energy}), retrieval entropy
(Equation~\eqref{eq:retrieval_entropy}), and normalised entropy
(Equation~\eqref{eq:normalised_entropy}) at position $(l,t)$, computed from
the retrieval scores of Equation~\eqref{eq:retrieval_scores}, and let
$\mathrm{JS}^{(l,t)}$ and $H^2_{\mathrm{Hel}}{}^{(l,t)}$ be the
Jensen--Shannon divergence and squared Hellinger distance between
$\{p_i^{(l,t)}\}$ and the reference distribution
$\{p_i^{(l,\mathrm{ref})}\}$.

All quantities are pooled over the full set of generated positions
$\{1, \ldots, T_g\}$. This choice is deliberate. The pooling rule consults
nothing beyond the generated sequence itself, so every feature is
computable at inference time without gold answers or labels.

For every layer $l$ the following five per-sample scalars are computed:
\begin{align}
    \Delta E^{(l)}
        &= \frac{1}{T_g} \sum_{t=1}^{T_g} E^{(l,t)}
           - E^{(l,\mathrm{ref})},
    \label{eq:delta_energy} \\[4pt]
    \Delta H^{(l)}
        &= \frac{1}{T_g} \sum_{t=1}^{T_g} H^{(l,t)}
           - H^{(l,\mathrm{ref})},
    \label{eq:delta_entropy} \\[4pt]
    \Delta \bar{H}^{(l)}
        &= \frac{1}{T_g} \sum_{t=1}^{T_g} \bar{H}^{(l,t)}
           - \bar{H}^{(l,\mathrm{ref})},
    \label{eq:delta_norm_entropy} \\[4pt]
    \overline{\mathrm{JS}}^{(l)}
        &= \frac{1}{T_g} \sum_{t=1}^{T_g} \mathrm{JS}^{(l,t)},
    \label{eq:mean_js} \\[4pt]
    \overline{H^2_{\mathrm{Hel}}}{}^{(l)}
        &= \frac{1}{T_g} \sum_{t=1}^{T_g} H^2_{\mathrm{Hel}}{}^{(l,t)}.
    \label{eq:mean_hellinger}
\end{align}
The mean over generated positions is the length-invariant reduction
required by Section~\ref{subsec:pa_aggregation}: it measures the typical
per-token shift and is insensitive to the number of generated tokens, so
the features cannot discriminate by answer length alone.

Each scalar measures a different aspect of how recognition-side retrieval
shifts between the question prefix and the answer generation.
$\Delta E^{(l)}$ captures the change in retrieval energy, mixing alignment
and query-norm components as analysed in
Section~\ref{subsec:pa_scale_confounds}; $\Delta \bar{H}^{(l)}$ captures the
change in retrieval concentration; and $\overline{\mathrm{JS}}^{(l)}$
together with $\overline{H^2_{\mathrm{Hel}}}{}^{(l)}$ measure the symmetric
distributional shift between the two phases from two geometrically distinct
angles.

Four of the five scalars enter the feature vector. Defining the per-group
layer vectors
\begin{equation}
    \mathbf{e}_q = \bigl[\Delta E^{(0)},\,\ldots,\,\Delta E^{(L-1)}\bigr]
                   \in \mathbb{R}^{L},
    \quad
    \mathbf{h}_q = \bigl[\Delta\bar{H}^{(0)},\,\ldots,\,
                   \Delta\bar{H}^{(L-1)}\bigr] \in \mathbb{R}^{L},
    \label{eq:group_vectors_1}
\end{equation}
\begin{equation}
    \mathbf{j}_q = \bigl[\overline{\mathrm{JS}}^{(0)},\,\ldots,\,
                   \overline{\mathrm{JS}}^{(L-1)}\bigr] \in \mathbb{R}^{L},
    \quad
    \mathbf{d}_q = \bigl[\overline{H^2_{\mathrm{Hel}}}{}^{(0)},\,\ldots,\,
                   \overline{H^2_{\mathrm{Hel}}}{}^{(L-1)}\bigr]
                   \in \mathbb{R}^{L},
    \label{eq:group_vectors_2}
\end{equation}
the feature vector is their concatenation:
\begin{equation}
    \boldsymbol{\phi}_q
    =
    \bigl[\,\mathbf{e}_q \;\|\; \mathbf{h}_q \;\|\; \mathbf{j}_q \;\|\;
    \mathbf{d}_q\,\bigr]
    \in \mathbb{R}^{4L}.
    \label{eq:feature_vector}
\end{equation}
The un-normalised entropy shift $\Delta H^{(l)}$ is the fifth scalar; it
is retained in the stored outputs but excluded from $\boldsymbol{\phi}_q$,
since $\Delta \bar{H}^{(l)}$ carries the same concentration signal in a
cross-layer-comparable form, as required by
Section~\ref{subsec:pa_distribution_comparison}. The absolute
generation-phase quantities $\overline{E}^{(l)}_{\mathrm{gen}}$ and
$\overline{\bar{H}}^{(l)}_{\mathrm{gen}}$, pooled without subtracting the
prefill reference, are likewise stored as diagnostics but excluded from the
feature vector. All $L$ layers enter $\boldsymbol{\phi}_q$, including early
layers whose scores may be dominated by representation magnitude rather
than discriminative structure; the per-feature standardisation of
Section~\ref{subsec:probe} places all features on a common scale, and the
L2 regularisation of the probe shrinks the coefficients of uninformative
ones.
% ----------------------------------------------------------------
\subsection{Inverse-Temperature Calibration}
\label{subsec:beta_calibration}
% ----------------------------------------------------------------

The inverse temperature $\beta$ controls the sharpness of the softmax over
the bank $\mathbf{M}^{(l)}$. When $\beta$ is too small the softmax is
near-uniform and the divergences vanish; when it is too large the softmax
collapses onto one row and the divergences saturate. Both extremes destroy
the discriminative signal, so $\beta$ is calibrated once per model against
a fixed target for the mean normalised entropy.

Calibration draws a warm-up subset of $N_c$ questions from the training
split and captures the retrieval scores of
Equation~\eqref{eq:retrieval_scores} at the prefill reference position of
every layer. For each candidate $\beta$ in a fixed, approximately
logarithmic grid, the
mean normalised entropy of the induced softmax is evaluated across all
captured layers and warm-up samples; this mean decreases monotonically in
$\beta$. The selected value $\beta^\star$ is the point at which the curve
crosses the operating target $\bar{H}^\star$, found by linear interpolation
in $\log \beta$ between the two bracketing grid points. If the curve does
not cross the target within the grid, $\beta^\star$ is clamped to the grid
endpoint whose mean normalised entropy is closest to $\bar{H}^\star$, and
the clamping is logged.

The target places the operating point inside the range of $\bar{H}$ over
which the divergences remain informative, away from both the uniform and
the collapsed extremes. The target and the surrounding band are operational
settings motivated by this saturation behaviour, not estimated quantities;
their concrete values, together with $N_c$ and the grid, are fixed in
Section~\ref{sec:experimental_config}, and the sensitivity of the results
to $\beta^\star$ is examined in Chapter~\ref{chap:results}. Restricting
calibration to the training split keeps test-set information out of
$\beta^\star$.

% ----------------------------------------------------------------
\subsection{Detection Probe and Evaluation Protocol}
\label{subsec:probe}
% ----------------------------------------------------------------

The detection model is the L2-regularised logistic-regression probe
motivated in Section~\ref{subsec:pa_aggregation}. With $\sigma(\cdot)$ the
logistic function, the probe estimates
\begin{equation}
    \hat{z}_q
    =
    \sigma\!\left(\boldsymbol{\phi}_q^{\!\top}\boldsymbol{w} + b\right),
    \qquad
    \boldsymbol{w} \in \mathbb{R}^{4L},\quad b \in \mathbb{R},
    \label{eq:logistic_probe}
\end{equation}
and is fit by minimising the regularised binary cross-entropy
\begin{equation}
    \min_{\boldsymbol{w},\,b}\;
    \frac{1}{N}\sum_{q=1}^{N}
    \mathrm{BCE}\!\left(z_q,\,\hat{z}_q\right)
    +
    \frac{1}{2C N}\,\|\boldsymbol{w}\|_2^2,
    \label{eq:logistic_objective}
\end{equation}
where $C > 0$ is the inverse-regularisation strength; this parameterisation
matches the numerical solver used in
Section~\ref{sec:implementation}. A standardisation step (zero mean, unit
variance, fit on the training fold only) precedes the logistic regression,
satisfying the common-scale requirement of
Section~\ref{subsec:pa_aggregation}, and per-column median imputation
handles the missing values that arise when activation extraction fails at a
single layer. Each coefficient in $\boldsymbol{w}$ maps to one layer-metric
pair, so every claim about where the signal lives is interpretable per
layer.

The $N$ labelled samples are partitioned, stratified on $z_q$ with a fixed
seed, into train, validation, and test splits. The test split is locked
before any probe is fit and is not inspected until the final step. The
hyperparameter $C$ is selected by stratified cross-validation within the
training split; the validation split takes no part in this selection and
enters only the final refit, a deliberately conservative use that keeps
model selection confined to the training data. An outer stratified
cross-validation over train~$\cup$~validation is used solely to report the
stability of the AUROC across folds, with the selected $C$ held fixed; it
plays no role in model selection and the resulting estimate is descriptive.
With the selected $C$ fixed, the probe is re-fit on
train~$\cup$~validation and evaluated once on the held-out test split, and
that single evaluation is the reported generalisation estimate. Split
ratios, fold counts, and the grid for $C$ are fixed in
Section~\ref{sec:experimental_config}.

Two diagnostics accompany the probe. The \emph{per-layer AUROC} of each
individual feature
$f \in \{\Delta E, \Delta\bar{H}, \overline{\mathrm{JS}},
\overline{H^2_{\mathrm{Hel}}}\}$ at every layer $l$ localises the signal
along the depth axis and feeds the exploratory depth-axis analysis of
Section~\ref{subsec:an_localisation}. The
\emph{single-feature unsupervised AUROC} of the best-layer score for each
$f$, computed without any logistic fitting, separates the unsupervised
contribution of the Hopfield features from the supervised contribution of
the probe.

% ----------------------------------------------------------------
\subsection{Algorithmic Summary}
\label{subsec:algorithms}
% ----------------------------------------------------------------

The banks are constructed once per model before any sample is processed:
for every layer, the gate matrix is extracted from the frozen weights as in
Equation~\eqref{eq:memory_bank_definition}, its rows are L2-normalised
(Equation~\eqref{eq:retrieval_scores}), and the per-row norms are stored so that
retrieval scores can be recovered from the captured gate pre-activations by
the fixed rescaling of Equation~\eqref{eq:retrieval_scores}.
Algorithm~\ref{alg:per_sample} then summarises the per-sample procedure. It
is purely functional: it takes the frozen model, the pre-built banks, the
calibrated $\beta^\star$, and the question, and returns the feature vector
$\boldsymbol{\phi}_q$ together with the generated text $\hat{y}$. No
in-place modification of the model occurs at any stage.

\begin{algorithm}[H]
\caption{Per-sample feature extraction}
\label{alg:per_sample}
\begin{algorithmic}[1]
\Require frozen model $f_\theta$, banks
         $\{\mathbf{M}^{(l)}\}_{l=0}^{L-1}$ with unit-normalised rows
         $\hat{\mathbf{m}}_i^{(l)}$ and stored row norms,
         inverse temperature $\beta^\star$, question $q$
\State $p \gets$ the question $q$ wrapped in the fixed chat template
\State Run the prefill forward pass on $p$; at every layer $l$, capture the
       hidden state $\mathbf{r}^{(l,\mathrm{ref})}$ at the last prompt token
\State Compute the retrieval scores of Equation~\eqref{eq:retrieval_scores},
       then $E^{(l,\mathrm{ref})}$, $H^{(l,\mathrm{ref})}$,
       $\bar{H}^{(l,\mathrm{ref})}$, and the reference distribution
       $\{p_i^{(l,\mathrm{ref})}\}$ for all $l$ via
       Equations~\eqref{eq:modern_hopfield_energy},
       \eqref{eq:retrieval_entropy}, \eqref{eq:normalised_entropy}, and
       \eqref{eq:retrieval_distribution}
\State Greedy-decode $\hat{y}$ from $p$ with $T_g \le T_{\max}$ tokens; at
       every generated position $t$, capture
       $\{\mathbf{r}_t^{(l)}\}_{l=0}^{L-1}$
\For{each $(l, t)$}
    \State Compute the retrieval scores, then $E^{(l,t)}$, $H^{(l,t)}$,
           $\bar{H}^{(l,t)}$, and $\{p_i^{(l,t)}\}$
    \State Compute $\mathrm{JS}^{(l,t)}$ and
           $H^2_{\mathrm{Hel}}{}^{(l,t)}$ against
           $\{p_i^{(l,\mathrm{ref})}\}$
\EndFor
\State Pool the per-token quantities over $t = 1, \ldots, T_g$ as in
       Equations~\eqref{eq:delta_energy}--\eqref{eq:mean_hellinger}
\State Assemble $\boldsymbol{\phi}_q$ as in
       Equations~\eqref{eq:group_vectors_1}--\eqref{eq:feature_vector}
\State \Return $(\boldsymbol{\phi}_q,\,\hat{y})$
\end{algorithmic}
\end{algorithm}
% =============================================================
% =============================================================
\section{Implementation}
\label{sec:implementation}
% =============================================================

The method is implemented as a Python package, \texttt{energy\_llm}, built
around an artefact pipeline. Each stage writes typed artefacts to disk and
reads only the artefacts produced by earlier stages, so every stage is
independently re-runnable and the expensive GPU stages are never repeated
when a downstream decision changes. The package is built on PyTorch
\citep{paszke2019pytorch} for tensor computation and activation capture,
the HuggingFace \texttt{transformers} library \citep{wolf2020transformers}
for model loading and tokenisation, NumPy and SciPy for the
information-theoretic primitives, and scikit-learn
\citep{pedregosa2011scikit} for the probe and evaluation protocol. Each
stage carries unit tests on its artefact contracts, and an end-to-end
integration test exercises the full pipeline on a miniature model with
\texttt{pytest}. The
specific classes and functions drawn from each library, and the reasons for
those choices, are stated where they are used.

% ----------------------------------------------------------------
\subsection{Stage 1: Bank Extraction}
\label{subsec:stage1}
% ----------------------------------------------------------------

Models are loaded through
\texttt{transformers.AutoModelForCausalLM.from\_pretrained}, which resolves
the architecture from the checkpoint configuration and reproduces the exact
parameterisation of the released weights; this matters because the bank is
defined as a function of those frozen weights and any re-implementation of
the architecture would break the identification of
Equation~\eqref{eq:memory_bank_definition}. The stage walks every decoder
layer, resolves the SwiGLU sub-modules through a model-alias registry
(\texttt{gate\_proj} and \texttt{up\_proj} for the Qwen and Llama
families), and reads the gate weight matrix. When the model is loaded with
\texttt{bitsandbytes} quantisation to fit GPU memory, the gate weights are
dequantised before bank construction, and in all cases they are cast to
\texttt{float32} on CPU; the row norms and the log-sum-exp of
Equation~\eqref{eq:modern_hopfield_energy} are sensitive to the reduced
mantissa of half-precision formats, and computing them in single precision
removes that error source at negligible cost, since the banks are built
once per model.

Row-wise L2 normalisation is then applied, and the per-row norms are stored
alongside the normalised rows, so that the raw gate pre-activations of
Equation~\eqref{eq:gate_preactivations} remain recoverable from the
retrieval scores by the fixed rescaling of
Equation~\eqref{eq:retrieval_scores}. The output artefact
\texttt{banks.pt}, serialised with \texttt{torch.save}, stores
$\{l \mapsto \texttt{Tensor}[K_l, D]\}$ together with a metadata record of
per-layer dimensions, the stored norms, and the bank identifier.

The primary and only evaluated configuration is the gate bank
$\mathbf{M}^{(l)} = \mathbf{A}^{(l)}$ of
Equation~\eqref{eq:memory_bank_definition}. The extraction interface also
supports the un-gated value branch $\mathbf{B}^{(l)}$ and the rows of
${\mathbf{W}_2}^{\!\top}$ as alternative configurations; these are
implementation capabilities rather than evaluated conditions, and no result
in this thesis depends on them. The raw \texttt{down\_proj} weight in its
native orientation is rejected explicitly, because its rows live in the
output dimension and would silently confuse keys with values.

% ----------------------------------------------------------------
\subsection{Stage 2: Trajectory Capture}
\label{subsec:stage2}
% ----------------------------------------------------------------

For each sample, this stage runs the two forward passes of
Section~\ref{subsec:pa_probe_positions} and realises
Algorithm~\ref{alg:per_sample}. The prompt is wrapped through the
tokenizer's \texttt{apply\_chat\_template} method, which applies the
template shipped with the checkpoint itself; using the shipped template
rather than a hand-written one guarantees that $q$ reaches the model in the
conversational format of its instruction tuning, and that prefill and
generation share the identical prompt prefix required by
Section~\ref{subsec:pa_probe_positions}. Decoding uses
\texttt{model.generate} with sampling disabled and a fixed
\texttt{max\_new\_tokens} cap, which implements the greedy argmax decoding
of Section~\ref{sec:problem_description} with $T_g \le T_{\max}$.

The probe point is captured with PyTorch forward pre-hooks
(\texttt{register\_forward\_pre\_hook}) registered on the \texttt{mlp}
module of every decoder layer. A pre-hook receives the exact input tensor
of the module it is attached to, which at this position is
$\mathbf{r}_t^{(l)}$, the output of the post-attention RMSNorm and the
argument of the gate and value projections
(Figure~\ref{fig:transformer_block}). The hook mechanism is chosen over
patching the model's forward method because it observes without rewriting:
the computation graph is untouched, which realises the read-only constraint
of this chapter mechanically rather than by convention. All passes run
under \texttt{torch.inference\_mode}, which disables gradient tracking and
makes any accidental parameter mutation an error. During prefill the hook
retains only the state at the last prompt position, yielding
$\mathbf{r}^{(l,\mathrm{ref})}$; during generation it retains the state of
each newly produced token.

Inside the hook, the captured state is multiplied by the row-normalised
bank, so the resulting vector is exactly the score vector
$\mathbf{s}_t^{(l)}$ of Equation~\eqref{eq:retrieval_scores}; the energy,
the retrieval distribution, and all derived quantities therefore operate on
the rescaled scores, never on the raw pre-activations, in accordance with
Section~\ref{subsec:retrieval_scoring}. The log-sum-exp term of the energy
is evaluated with \texttt{scipy.special.logsumexp}, which shifts by the
maximum score before exponentiating and thereby avoids the overflow that a
literal evaluation of Equation~\eqref{eq:modern_hopfield_energy} produces
at large $\beta$; the log-sum-exp and quadratic terms are stored as
separate fields, preserving the decomposition required by
Section~\ref{subsec:pa_scale_confounds}. The Jensen--Shannon divergence
requires a convention reconciliation:
\texttt{scipy.spatial.distance.jensenshannon} returns the square root of
the divergence in the logarithm base passed to it, so the package squares
the returned value and converts to natural logarithms,
\begin{equation}
    \mathrm{JS}_{\mathrm{nats}}(P, Q)
    =
    \bigl[\texttt{jensenshannon}(P, Q,\ \mathrm{base}=2)\bigr]^{2}
    \cdot \log 2,
    \label{eq:js_convention}
\end{equation}
which pins every reported divergence to the nats convention of
Section~\ref{sec:divergences} and keeps the bound
$\mathrm{JS} \le \log 2$ exact. The squared Hellinger distance is
implemented directly in NumPy from Equation~\eqref{eq:hellinger_squared},
since it involves only square roots and a sum and needs no library wrapper.

Hidden states themselves are never persisted: the bank multiplication and
the divergences are computed online inside the hook, and only the resulting
$[L]$ and $[L, T_g]$ scalar arrays are written to a compressed per-sample
artefact, together with a metadata file recording the question, the gold
answers, the generated text, and the token count. The stage exposes a
sharding interface that slices the dataset into non-overlapping chunks by
stride, so several workers can populate the same trajectory directory in
parallel, as used in the HPC array-job submission; the per-sample
aggregation of Stage~4 is unaffected by the order in which shards complete.

% ----------------------------------------------------------------
\subsection{Inverse-Temperature Calibration}
\label{subsec:impl_beta}
% ----------------------------------------------------------------

The calibration of Section~\ref{subsec:beta_calibration} is realised as a
warm-up pass that precedes the main Stage~2 run. The prefill reference
scores of $N_c$ training-split questions are captured once, and the mean
normalised entropy of the induced softmax is evaluated for every $\beta$ on
a fixed, approximately logarithmic grid by recomputing the softmax from the
cached score
vectors, without further forward passes. The selected $\beta^\star$ is
obtained by linear interpolation in $\log \beta$ between the two grid
points bracketing the target $\bar{H}^\star$; when the target is not
bracketed, the nearest grid endpoint is taken and the clamping is recorded
in the run metadata. The calibrated value is written into the experiment
configuration snapshot so that every subsequent stage reads the same
$\beta^\star$.

% ----------------------------------------------------------------
\subsection{Stage 3: Labelling}
\label{subsec:stage3_labeling}
% ----------------------------------------------------------------

Labels are produced by the external judge of
Section~\ref{sec:problem_description}, queried through an OpenAI-compatible
chat-completions interface with the locked prompt template that presents
$q$, $\hat{y}$, and $\mathcal{A}_q$ and requests a binary supported or
unsupported verdict. The raw judge response is archived verbatim alongside
the parsed label in a per-sample label artefact, so the parsing rule can be
audited and re-applied without re-querying the service. Labels are stored
separately from the trajectory artefacts, which makes relabelling a
CPU-only operation that never touches the GPU stages. The agreement of the
judge with a human-labelled subset, reported as Cohen's $\kappa$, is
computed in this stage and carried into Chapter~\ref{chap:results}.

% ----------------------------------------------------------------
\subsection{Stage 4: Feature Aggregation and Probing}
\label{subsec:stage4}
% ----------------------------------------------------------------

This stage is CPU-only and reads only the artefacts of Stages~2 and~3. For
each sample it loads the per-token scalar arrays, pools them over the full
set of generated positions $\{1, \ldots, T_g\}$ as specified in
Section~\ref{subsec:feature_design}, computes the per-layer features of
Equations~\eqref{eq:delta_energy}--\eqref{eq:mean_hellinger}, and assembles
the matrix whose rows are the vectors $\boldsymbol{\phi}_q$ of
Equation~\eqref{eq:feature_vector}. No gold answer enters this computation;
the pooling consults only the generated sequence, so the features remain
computable at inference time, and the labels of Stage~3 are joined to the
feature matrix only at fitting time.

The probe of Section~\ref{subsec:probe} is realised as a scikit-learn
\texttt{Pipeline} chaining \texttt{SimpleImputer} with the median strategy,
\texttt{StandardScaler}, and \texttt{LogisticRegression} with the
\texttt{lbfgs} solver. The \texttt{Pipeline} class is used deliberately
rather than applying the preprocessing steps manually: when the pipeline is
passed to the cross-validation utilities, the imputer and the scaler are
re-fit inside every training fold and applied unchanged to the
corresponding held-out fold, which enforces by construction that no
test-fold statistics leak into the standardisation. Splitting and fold
generation use \texttt{StratifiedKFold} and a stratified
\texttt{train\_test\_split} with fixed seeds, preserving the class balance
of $z_q$ in every partition, and discrimination is computed with
\texttt{roc\_auc\_score}. The scikit-learn objective is parameterised as
$\tfrac{1}{2}\|\mathbf{w}\|_2^2 + C\sum_q \mathrm{BCE}$, which coincides
with Equation~\eqref{eq:logistic_objective} up to the positive factor
$CN$, leaving the minimiser unchanged; the
$C$ values reported in this work refer to the scikit-learn convention. The
per-layer AUROC and single-feature diagnostics of
Section~\ref{subsec:probe} are computed in the same stage directly from the
feature matrix, with no fitted model involved.

% =============================================================
% =============================================================
% =============================================================
% =============================================================
\section{Analysis Method}
\label{sec:analysis_method}
% =============================================================

This section fixes the statistical procedures behind the quantitative
claims of Chapter~\ref{chap:results}. The thesis makes four claims that
require statistical support: that the probe discriminates hallucinated from
grounded answers, that the Hopfield features contribute information beyond
what the generation process already exposes, that the signal admits a
depth-axis localisation consistent with the mechanistic objective of this
chapter, and that the judge labels supporting all of the above are
reliable. One procedure is fixed per claim, and the section closes by
separating the confirmatory claims from the exploratory ones.

% ----------------------------------------------------------------
\subsection{Discrimination of the Probe}
\label{subsec:an_discrimination}
% ----------------------------------------------------------------

The primary endpoint is the AUROC of the probe on the locked test split of
Section~\ref{subsec:probe}, reported with a non-parametric bootstrap $95\%$
confidence interval from $B = 10\,000$ resamples of the test predictions;
the resample count controls only the numerical stability of the interval,
while the width of the interval reflects the size of the test split. A
lower confidence bound above $0.5$ is the condition for claiming
discrimination. Resamples containing a single class, for which the AUROC is
undefined, are redrawn. Because the test split is small relative to the category
structure of TruthfulQA, no stratified inferential procedure is attempted
at the category level; instead, the AUROC computed separately within the
largest categories is reported descriptively, to verify that the
discrimination is not the artefact of a single category with an atypical
hallucination rate.

% ----------------------------------------------------------------
\subsection{Contribution of the Hopfield Features}
\label{subsec:an_contribution}
% ----------------------------------------------------------------

Discrimination alone does not establish the thesis: a detector could
discriminate using only the model's own confidence in its output. The
central claim is therefore comparative. Two baseline scalars are computed
from the same greedy decoding pass: the length-normalised sequence
log-probability and the mean token-level predictive entropy of the
generated answer. Two probes are fit under the identical protocol of
Section~\ref{subsec:probe}, a base probe on the two baselines alone and a
full probe on the baselines concatenated with $\boldsymbol{\phi}_q$. For
each of $B = 10\,000$ paired bootstrap resamples of the test set,
$\Delta\mathrm{AUROC} = \mathrm{AUROC}^{\mathrm{full}} -
\mathrm{AUROC}^{\mathrm{base}}$ is recomputed, and the Hopfield
contribution is claimed significant only if the lower bound of the
resulting $95\%$ confidence interval is positive. The pairing is essential:
both probes are evaluated on the same resample, so the variance of the
difference excludes the sampling noise common to both and the test isolates
the added value of the features.

% ----------------------------------------------------------------
\subsection{Localisation Along the Depth Axis}
\label{subsec:an_localisation}
% ----------------------------------------------------------------

The per-layer AUROC profiles of Section~\ref{subsec:probe} serve the
interpretability objective and are exploratory; they support no
confirmatory claim. Because the direction of the effect is unknown per
layer (Section~\ref{subsec:pa_aggregation}), the per-layer statistic is the
folded value $\max(\mathrm{AUROC},\,1-\mathrm{AUROC})$. The best layer is
then a scan statistic: even under pure chance, the maximum over $L$ layers
exceeds any pointwise threshold with high probability, so the profile is
protected by a permutation null of the scan maximum. Labels are shuffled,
the maximum folded AUROC over all layer and feature combinations entering
the observed profiles is recomputed, and the procedure is
repeated $B_{\mathrm{perm}} = 1000$ times; a layer band is described as
carrying signal only where the observed values exceed the $95$th percentile
of this null.

% ----------------------------------------------------------------
\subsection{Label Reliability}
\label{subsec:an_label_reliability}
% ----------------------------------------------------------------

Every claim above inherits the quality of the judge labels of
Section~\ref{sec:problem_description}. Agreement between the judge and the
human-labelled subset is quantified by Cohen's $\kappa$, reported together
with the subset size. A value of $\kappa \ge 0.6$, substantial agreement
under the convention of \citet{landis1977measurement}, is the condition for
treating the judge labels as an acceptable ground truth; if the threshold
is not met, the affected results are reported with this limitation stated
explicitly rather than withheld.

% ----------------------------------------------------------------
\subsection{Sensitivity and Confirmatory Family}
\label{subsec:an_sensitivity}
% ----------------------------------------------------------------

The sensitivity of the test AUROC to the calibrated inverse temperature is
examined descriptively by recomputing the features over the $\beta$ grid of
Section~\ref{subsec:beta_calibration} from the cached score tensors,
without further forward passes, and reporting the AUROC as a function of
$\beta$. The confirmatory family consists of exactly two hypotheses, the
discrimination claim of Section~\ref{subsec:an_discrimination} and the
contribution claim of Section~\ref{subsec:an_contribution}, tested with
Bonferroni correction at family level $\alpha = 0.05$. All remaining
analyses, including the depth-axis profiles, the per-category breakdown,
and the $\beta$ sensitivity curve, are exploratory or descriptive and carry
no corrected significance claim.

% =============================================================
% =============================================================
\section{Experimental Configuration}
\label{sec:experimental_config}
% =============================================================

This section fixes the configuration under which the final pipeline runs.
Nothing here overrides the design decisions above; the purpose is to set
every free parameter before execution so that the results of
Chapter~\ref{chap:results} are reproducible from a single configuration
file.

% ----------------------------------------------------------------
\subsection{Models}
\label{subsec:models}
% ----------------------------------------------------------------

Two instruction-tuned decoder-only models are evaluated.
\textbf{Qwen2.5-3B-Instruct} \citep{qwen2025qwen25}
($L = 36$, $D = 2048$, $D_{\mathrm{ff}} = 11008$) is the primary model,
used for calibration, hyperparameter selection, and the headline result.
\textbf{Llama-3.2-3B-Instruct} \citep{grattafiori2024llama3}
($L = 28$, $D = 3072$, $D_{\mathrm{ff}} = 8192$) is the replication model
for cross-architecture transfer. All reported runs execute on a single
A100 GPU in BF16; 4-bit NF4 quantisation (\texttt{bitsandbytes}) is used
only for local development and produces no reported number, in keeping
with the fixed-precision determinism constraint of this chapter. Bank
extraction always dequantises the gate matrix to \texttt{float32} on CPU
before row-wise L2 normalisation
(Section~\ref{subsec:stage1}).

% ----------------------------------------------------------------
\subsection{Datasets}
\label{subsec:datasets}
% ----------------------------------------------------------------

\textbf{TruthfulQA} \citep{lin2022truthfulqa} ($817$ adversarial questions
across $38$ categories) is the primary dataset: it provides curated
correct and incorrect reference answers, is a standard benchmark for
hallucination detection on factoid QA, and its category metadata supports
the descriptive per-category breakdown of
Section~\ref{subsec:an_discrimination}. \textbf{TriviaQA}
(\texttt{rc.nocontext} split, \citealt{joshi2017triviaqa}, $500$ samples
drawn at the fixed seed) provides a secondary cross-dataset evaluation,
executed only if the primary analysis completes within the compute budget;
its results are reported as exploratory in either case.

% ----------------------------------------------------------------
\subsection{Experiments}
\label{subsec:experiments}
% ----------------------------------------------------------------

Three runs are defined, in strict priority order. \textbf{E1} (primary):
Qwen2.5-3B-Instruct on TruthfulQA; this run supports both confirmatory
hypotheses of Section~\ref{subsec:an_sensitivity}, the depth-axis profile,
the per-category breakdown, and the $\beta$ sensitivity curve. \textbf{E2}
(replication): Llama-3.2-3B-Instruct on TruthfulQA under the identical
configuration, with $\beta^\star$ recalibrated for the model; this run
tests whether the discrimination and contribution findings transfer across
architectures and is reported descriptively against the E1 outcome.
\textbf{E3} (secondary, contingent): Qwen2.5-3B-Instruct on TriviaQA,
reported as exploratory cross-dataset evidence. The $\beta$ sensitivity
analysis requires no additional GPU run, since it recomputes features from
the score tensors cached during E1
(Section~\ref{subsec:an_sensitivity}). For the label-reliability check of
Section~\ref{subsec:an_label_reliability}, a random subset of $60$ E1
samples is labelled by hand and compared against the judge.

% ----------------------------------------------------------------
\subsection{Hyperparameters}
\label{subsec:hyperparameters}
% ----------------------------------------------------------------

Table~\ref{tab:hyperparameters} lists every hyperparameter of the final
pipeline. These values are fixed before execution and are not varied
across the runs that produce the headline result.

\begin{table}[H]
    \centering
    \small
    \begin{tabular}{p{0.36\linewidth} p{0.28\linewidth} p{0.28\linewidth}}
        \toprule
        Parameter & Value & Provenance \\
        \midrule
        \multicolumn{3}{l}{\emph{Memory bank}} \\
        \hspace{1em}\texttt{bank}
            & \texttt{gate}
            & Eq.~\eqref{eq:memory_bank_definition} \\
        \hspace{1em}\texttt{normalize}
            & \texttt{True} (row L2)
            & \S\ref{subsec:retrieval_scoring} \\
        \hspace{1em}\texttt{hook\_target}
            & \texttt{mlp\_input}
            & \S\ref{subsec:pa_probe_positions} \\
        \multicolumn{3}{l}{\emph{Energy and divergences}} \\
        \hspace{1em}\texttt{score}
            & gate pre-activation, rescaled
            & Eq.~\eqref{eq:retrieval_scores} \\
        \hspace{1em}\texttt{normalize\_query}
            & \texttt{False}
            & \S\ref{subsec:retrieval_scoring} \\
        \hspace{1em}$\beta^\star$
            & calibrated, target $\bar{H}^\star{=}0.55$
            & \S\ref{subsec:beta_calibration} \\
        \hspace{1em}$\beta$ grid
            & $\{1,2,5,10,15,20,30,50,100\}$
            & \S\ref{subsec:beta_calibration} \\
        \hspace{1em}Operating band
            & $\bar{H} \in [0.3,\,0.7]$
            & \S\ref{subsec:beta_calibration} \\
        \hspace{1em}\texttt{calibration\_samples}
            & $N_c = 30$
            & \S\ref{subsec:beta_calibration} \\
        \multicolumn{3}{l}{\emph{Generation}} \\
        \hspace{1em}Decoding
            & greedy
            & \S\ref{sec:problem_description} \\
        \hspace{1em}$T_{\max}$
            & $50$ tokens
            & \S\ref{subsec:feature_design} \\
        \hspace{1em}Chat template
            & checkpoint default
            & \S\ref{subsec:stage2} \\
        \multicolumn{3}{l}{\emph{Feature aggregation}} \\
        \hspace{1em}Pooling
            & mean over all generated tokens
            & Eqs.~\eqref{eq:delta_energy}--\eqref{eq:mean_hellinger} \\
        \multicolumn{3}{l}{\emph{Baselines}} \\
        \hspace{1em}Baseline features
            & seq.\ log-prob.\ (length-norm.), mean token entropy
            & \S\ref{subsec:an_contribution} \\
        \multicolumn{3}{l}{\emph{Probe}} \\
        \hspace{1em}Regulariser
            & L2
            & Eq.~\eqref{eq:logistic_objective} \\
        \hspace{1em}$C$ grid
            & $\{0.01,\,0.1,\,1,\,10\}$
            & \S\ref{subsec:probe} \\
        \hspace{1em}Inner CV
            & $3$-fold stratified
            & \S\ref{subsec:probe} \\
        \hspace{1em}Outer CV
            & $5$-fold stratified
            & \S\ref{subsec:probe} \\
        \hspace{1em}Split
            & $70/15/15$
            & \S\ref{subsec:probe} \\
        \multicolumn{3}{l}{\emph{Statistics}} \\
        \hspace{1em}Bootstrap resamples
            & $B = 10\,000$
            & \S\ref{subsec:an_discrimination} \\
        \hspace{1em}Permutation resamples
            & $B_{\mathrm{perm}} = 1000$
            & \S\ref{subsec:an_localisation} \\
        \hspace{1em}Confirmatory family
            & Bonferroni, $\alpha{=}0.05$, $2$ tests
            & \S\ref{subsec:an_sensitivity} \\
        \multicolumn{3}{l}{\emph{Labelling}} \\
        \hspace{1em}Judge model
            & \texttt{gpt-4o-mini}
            & \S\ref{sec:problem_description} \\
        \hspace{1em}Human-labelled subset
            & $60$ samples
            & \S\ref{subsec:an_label_reliability} \\
        \multicolumn{3}{l}{\emph{Determinism}} \\
        \hspace{1em}Random seed
            & $42$
            & --- \\
        \bottomrule
    \end{tabular}
    \caption{Final pipeline hyperparameters. The provenance column gives
    the section in which each choice is justified.}
    \label{tab:hyperparameters}
\end{table}

% =============================================================
\section{Summary}
\label{sec:method_summary}
% =============================================================

The method turns the theoretical framework of Chapter~\ref{chap:theory} into a
concrete, auditable pipeline. The probe treats the hidden state entering each
SwiGLU sublayer as a query against the layer's gate matrix $\mathbf{A}^{(l)}$,
whose rows are the pattern detectors of the feed-forward sublayer. The
retrieval scores are a fixed per-row rescaling of the gate pre-activations the
model already computes; as Section~\ref{subsec:retrieval_scoring} makes
explicit, they are a deterministic derived view of the matching step rather
than the raw activations themselves. The method computes the retrieval energy
and the divergences against the prefill last-token reference, pools the
per-token quantities over all generated tokens, and feeds the per-layer
features into an L2-regularised logistic regression evaluated on a locked
$15\%$ test split. The gate bank is the primary and only evaluated
configuration; the value-branch and value-space banks are retained as
implementation capabilities (Section~\ref{subsec:stage1}) and support no
result reported in this thesis. The four-stage implementation writes typed
artefacts at every step, supports parallel sharded execution on HPC
infrastructure, and is validated by unit and integration tests. The
configuration of Section~\ref{sec:experimental_config} is the one that
Chapter~\ref{chap:results} executes and reports against pre-registered
baselines and statistical procedures.